% Source observation: analyses/018-2026-09-08-results-materials/observations/O007-kernel-realization.md
% Replacement prose for kernel-autoresearch.tex; Analysis 018 O007 and figure_data.json:runtime.
% Source reduction: evidence.py:runtime_subset; figure: plots.py:kernels.
% Underlying implementation audit: Run 029 observations/06-implementation-and-claim-audit.md.
% Argument revision: Analysis 018 observations/O015-threshold-kernel-argument.md.
% Compatibility evidence: Run 022 README; Run 025 SOURCE_AUDIT.md; Run 029 Observation 04.
\subsection{Realizing sparsity with specialized inference kernels}
\label{sec:kernel-autoresearch}

Realizing activation sparsity requires kernels matched to the model's shapes,
nonlinearities and hardware. The released TwELL kernels~\citep{cetin2026sparser}
accelerated their supported workload, but the unchanged small-width primitive
failed numerical qualification for Pythia-14M. A shape- and sign-adapted
descendant, P0, qualifies on only four of our 30 checkpoints, with geometric
mean speedup $0.8910\times$ on that subset (Appendix~\ref{app:kernel-results}).
This motivates specialization for the architecture and intervention.

To make this specialization practical within an architecture study, we use
a human-guided coding-agent loop to propose, test and refine implementations.
Execution and profiling feedback provide a basis for this
approach~\citep{ouyang2025kernelbench,luo2026agentic}.
We replay 42 eligible proposals from the Codex search against native PyTorch
scaled dot-product attention (SDPA) in CUDA graphs on one RTX5090.
The workload is BF16, batch-one, uncached 2,048-token inference with full
vocabulary logits. Qualification requires logits and loss to agree with the
native reference within specified tolerances on complete validation.

The matched retrospective reaches $1.6024\times$ by iteration 11 and
$1.7830\times$ at iteration 42 on the fixed A7-OL1, $\kappa=0.5$ search
checkpoint (Figure~\ref{fig:kernel-autoresearch}, left): substantial gains
appear early, followed by a later improvement.
The selected implementation K050 qualifies on all 30 checkpoints in the
14M cohort, with geometric mean speedup $1.2340\times$. Speedup has a strong,
approximately linear association with model-wide sparsity
(right; OLS $R^2=0.8167$). This connects the logical metric to measured
execution, while checkpoint quality, topology and weights vary across the fit.

Ablations distinguish fusion from the contribution of sparse paths.
Relative to the same fusion with sparse paths disabled, K050 is
$1.311\times$ as fast at the search checkpoint and $1.043\times$ as fast
geometrically across the cohort, helping 14 of 30 checkpoints.
Attention skipping is unprofitable in this workload:
disabling it improves every checkpoint, raising the mean to $1.2506\times$.
The attention kernel can skip a matrix-multiply instruction when either
complete operand fragment is zero; matching zeros in both $Q$ and $K$ are
unnecessary. Yet even substantial instruction skipping fails to repay its
overhead at 14M (Appendix~\ref{app:kernel-results}). The full-model gain thus
combines fusion with a conditional benefit from sparse paths: profitable
execution depends on both sparsity and the granularity at which work is skipped.

\begin{figure}[!htbp]
\centering
\includegraphics[width=\linewidth]{figures/07-kernel-realization.pdf}
\caption{\textbf{Sparse kernel auto-research for full-model acceleration.}
(a) Matched retrospective of the best qualified incumbent on fixed A7-OL1, $\kappa=0.5$, across 42
iterations, starting at native $1\times$ and reaching $1.7830\times$.
(b) The selected K050 implementation on 30 Pythia-14M checkpoints; symbols identify recipes. The dashed
line is a descriptive OLS fit with intercept ($R^2=0.8167$); dotted lines
mark native $1\times$. Speedup uses each checkpoint's native SDPA CUDA-graph
baseline on one RTX5090: BF16, batch one, uncached 2,048-token inference and
full vocabulary logits. Qualification covers all 338 validation blocks. All checkpoints
qualify, with geometric mean $1.2340\times$. Here ``best'' refers to the
searched proposals; the attention-dense ablation is faster on every checkpoint
(Table~\ref{tab:kernel-ablations}). Sparsity uses FP16 measurements, while
timing uses BF16; both speedup axes have restricted ranges. The fit describes
association across different checkpoints, not a causal sparsity effect.}
\label{fig:kernel-autoresearch}
\end{figure}
\FloatBarrier
